% Standalone problem document, generated from the adjacent README.md.
% Compile with XeLaTeX. Mathematical content is maintained in Markdown.
\documentclass[11pt,a4paper]{article}
\usepackage[fontsize=10.5pt]{scrextend}
\usepackage[margin=22mm,headheight=15pt,headsep=9mm,footskip=11mm]{geometry}
\usepackage{fontspec}
\setmainfont{texgyrepagella}[Extension=.otf,UprightFont=*-regular,BoldFont=*-bold,ItalicFont=*-italic,BoldItalicFont=*-bolditalic]
\setsansfont{texgyreheros}[Extension=.otf,UprightFont=*-regular,BoldFont=*-bold,ItalicFont=*-italic,BoldItalicFont=*-bolditalic]
\setmonofont{texgyrecursor}[Extension=.otf,UprightFont=*-regular,BoldFont=*-bold,ItalicFont=*-italic,BoldItalicFont=*-bolditalic,Scale=MatchLowercase]
\usepackage{amsmath,amssymb}
\usepackage{unicode-math}
\setmathfont{texgyrepagella-math.otf}
\usepackage{xcolor,microtype,enumitem,needspace,fancyhdr}
\usepackage{bookmark}
\definecolor{ink}{HTML}{17344B}
\definecolor{accent}{HTML}{197D80}
\definecolor{muted}{HTML}{596773}
\hypersetup{colorlinks=true,linkcolor=accent,urlcolor=accent,pdftitle={RA-02 - Polynomial
trace-error factor after exactly the target rank of
pivots},pdfauthor={Open Problems in Numerical Linear Algebra}}
\urlstyle{same}
\setlength{\parindent}{0pt}
\setlength{\parskip}{5pt plus 1pt minus 1pt}
\linespread{1.04}
\setlength{\emergencystretch}{3em}
\widowpenalty=10000
\clubpenalty=10000
\displaywidowpenalty=10000
\allowdisplaybreaks[1]
\setlist{nosep,leftmargin=1.5em}
\providecommand{\tightlist}{\setlength{\itemsep}{0pt}\setlength{\parskip}{0pt}}
\providecommand{\pandocbounded}[1]{#1}
\pagestyle{fancy}
\fancyhf{}
\fancyhead[L]{\sffamily\footnotesize\color{muted}OPEN PROBLEMS IN NUMERICAL LINEAR ALGEBRA}
\fancyhead[R]{\sffamily\bfseries\footnotesize\color{accent}RA-02}
\fancyfoot[L]{\parbox[b]{0.9\textwidth}{\sffamily\fontsize{8}{10}\selectfont\color{muted}Editorial ratings; a bounded search cannot certify that a problem remains unsolved.\\Literature check: 2026-09-11}}
\fancyfoot[R]{\sffamily\footnotesize\color{muted}\thepage}
\renewcommand{\headrulewidth}{0.3pt}
\renewcommand{\headrule}{\hbox to\headwidth{\color{accent}\leaders\hrule height \headrulewidth\hfill}}
\makeatletter
\renewcommand{\section}{\Needspace{5\baselineskip}\@startsection{section}{1}{0pt}{12pt plus 2pt minus 2pt}{4pt}{\sffamily\bfseries\large\color{ink}}}
\renewcommand{\subsection}{\Needspace{4\baselineskip}\@startsection{subsection}{2}{0pt}{9pt plus 1pt minus 1pt}{3pt}{\sffamily\bfseries\normalsize\color{ink}}}
\makeatother
\setcounter{secnumdepth}{0}
\begin{document}
{\sffamily\small\color{muted}Randomized and low-rank approximation\par}
\vspace{3pt}
{\raggedright\hyphenpenalty=10000\exhyphenpenalty=10000\sffamily\bfseries\fontsize{21}{24}\selectfont\color{ink}Polynomial
trace-error factor after exactly the target rank of pivots\par}
\vspace{7pt}
{\sffamily\small\textbf{Difficulty:} challenging\quad\textbf{Importance:} interesting
to the community\par}
{\sffamily\small\bfseries\color{accent}Status: Solved\par}
\vspace{4pt}
{\color{accent}\hrule height 0.8pt}
\vspace{6pt}
\textbf{Rating rationale:} Challenging because removing exponential loss
without oversampling requires sharper adaptive-pivot analysis; community
impact is rank-efficient PSD approximation.\\
\textbf{Topic:} randomized factorization; approximation guarantees

\subsection{Resolution --- 2026-09-11}\label{resolution-2026-09-11}

\textbf{Negative resolution by Matthew J. Colbrook}, Department of
Applied Mathematics and Theoretical Physics, University of Cambridge.
Theorem 1 and Corollary 3 disprove the existence of constants \(C,p\)
giving the displayed polynomial bound after exactly \(r\) pivots. For
every fixed \(r\ge1\), real entrywise-positive positive-definite
matrices of order \(r+1\) approach the sharp expected trace-error ratio
\(2^r\) as a parameter tends to zero. Choose \(r\) first and then the
parameter; no limit uniform in \(r\) is needed. The result does not
address oversampling (RA-01).

\href{https://github.com/ajt60gaibb/OpenProblemsInNLA/blob/main/references/colbrook-random-pivoting-2026-09-11/manuscripts/sharp_random_pivoting.pdf}{Complete
manuscript} ·
\href{https://github.com/ajt60gaibb/OpenProblemsInNLA/blob/main/references/colbrook-random-pivoting-2026-09-11/manuscripts/sharp_random_pivoting.tex}{TeX}
·
\href{https://github.com/ajt60gaibb/OpenProblemsInNLA/blob/main/references/colbrook-random-pivoting-2026-09-11/verification/reviews/RA-02-review.md}{Independent
complete-source PASS review} ·
\href{https://github.com/ajt60gaibb/OpenProblemsInNLA/blob/main/references/colbrook-random-pivoting-2026-09-11/README.md}{Authorship,
exact checks and provenance}.

Verification is independent agent review, not external human peer review
or formal certification. AI assistance is disclosed; no priority claim
is made. The original statement and audits remain below; ratings are
historical.

\subsection{Context and notation}\label{context-and-notation}

For a Hermitian positive-semidefinite \(A\in\mathbb C^{n\times n}\),
define the exact-arithmetic RPCholesky residuals by \(R_0=A\).
Conditional on \(R_t\ne0\), select \(j\) with probability
\((R_t)_{jj}/\operatorname{tr}(R_t)\) and set

\[
R_{t+1}=R_t-\frac{R_t(:,j)R_t(j,:)}{(R_t)_{jj}}.
\]

If \(R_t=0\), keep all subsequent residuals zero. Eigenvalues are
ordered \(\lambda_1(A)\geq\cdots\geq\lambda_n(A)\geq0\), and
\(\tau_r(A)=\sum_{j>r}\lambda_j(A)\). This is the pivot rule in Chen,
Epperly, Tropp, and Webber,
\href{https://doi.org/10.1002/cpa.22234}{\emph{Randomly pivoted
Cholesky: Practical approximation of a kernel matrix with few entry
evaluations}}, Algorithm 1; their Lemma 5.5 gives the current comparison
\(\mathbb E\operatorname{tr}(R_r)\leq2^r\tau_r(A)\).

\subsection{Problem statement}\label{problem-statement}

Do constants \(C>0\) and \(p\geq0\) exist such that, for every
\(n\geq1\), every Hermitian positive-semidefinite
\(A\in\mathbb C^{n\times n}\), and every integer \(1\leq r\leq n\),

\[
\mathbb E\operatorname{tr}(R_r)\leq Cr^p\tau_r(A)?
\]

Both constants must be independent of \(n,r,A\). Exactly \(r\)
RPCholesky steps are permitted, with the zero-residual convention above.
This asks for a polynomial approximation factor without oversampling;
\href{https://github.com/ajt60gaibb/OpenProblemsInNLA/blob/main/randomized-and-low-rank-approximation/RA-01/README.md}{RA-01}
instead allows additional pivots to obtain relative error near one.

\newpage
\subsection{References}\label{references}

Epperly,
\href{https://tropp.caltech.edu/dissertations/Epp25-Making-Most.pdf}{\emph{Make
the Most of What You Have}}, §11.1, Conjecture 11.2, p.~182; Chen et
al., \href{https://doi.org/10.1002/cpa.22234}{RPCholesky}, Lemma 5.5,
p.~1020. Gilles and Wilber,
\href{https://arxiv.org/html/2601.22344v1}{\emph{Low-Rank Approximation
by Randomly Pivoted LU}}, §3.1, paragraph following Theorem 3,
explicitly reiterate this conjecture.

\subsection{Source normalization}\label{source-normalization}

The dissertation calls \(A^{(r)}\) a residual but prints
\(\operatorname{tr}(A-A^{(r)})\), without expectation. We use its cited
Lemma 5.5 to correct both notation defects. The later paper confirms the
intended polynomial improvement over that lemma's \(2^r\) factor.

\subsection{Status check}\label{status-check}

Searches for \textit{RPCholesky r-step polynomial},
\textit{RPCholesky conjecture polynomial}, and the cited papers found
no resolution. Epperly's August 2026 oversampling theorem does not
supply a polynomial factor at exactly \(r\) steps. Its Theorem 1.2 is
noninformative at \(k=r\).

\subsection{Audit --- 2026-09-10}\label{audit-2026-09-10}

Rechecked \href{https://arxiv.org/html/2601.22344v1}{Gilles--Wilber's
discussion after Theorem 3}, which reiterates the polynomial-factor
conjecture. Exactly-\(r\)-step and later RPCholesky searches found no
resolution. The \href{https://arxiv.org/html/2608.20633v1}{August
oversampling theorem} does not establish the displayed factor after
exactly \(r\) pivots.
\end{document}
